% w_captions.tex -- GENERATED by figs_wnet.py (do not edit by hand; rerun the script).
% Caption macro for figs/w_net.pdf; every number is asserted by the script.
\newcommand{\wNetCaption}{%
Locked and free wraps at $m=3$, from points generated by the generator matrices, in cell units. Grey segments join the points of consecutive cells along the vertical axis. (a) $\mathcal T_{101}$ (Proposition~\ref{prop:twocoset}), a slab $n_z=\mathrm{const}$ seen in the $(y,x)$-plane, with the free coordinate $x$ vertical (shaded: cells raised by about $\frac12$ in $x$): the step along $\e_x$ moves $x$ by exactly one cell, up to the jitter, whatever the cell. (b) $\mathcal W_{101}$ (Proposition~\ref{prop:wnet}, $\varepsilon=\rho/(2p)=5/202$), a slab $n_y=\mathrm{const}$ in the $(x,z)$-plane, with the locked coordinate $z$ vertical: the step along $\e_z$ is $(4\varepsilon,4\varepsilon,1+2\varepsilon)$, longer than one cell (accent arrow); where $z$ wraps, $x$ and $y$ wrap too and the step becomes $(-1+4\varepsilon,-1+4\varepsilon,2\varepsilon)$, still long (thick accent arrow; checked at every cell of the window). (c) The same net after the digital shift that adds $1$ to the second digit of $z$: the wraps of $z$ are no longer locked, $z$ now wraps alone at some cells (the cells with $\zeta=p-1-\rho$), and the step along $\e_z$ there shrinks to $(4\varepsilon,4\varepsilon,2\varepsilon)$ up to the third digits (circles; length $0.148$); the shifted net has $\delta N^{1/3}=.1355\ldots$, against $1.0362\ldots$ without the shift (computed, Remark~\ref{rem:wshift}). (d) $\delta N^{1/3}$ of $\mathcal W_p$ (exact at the primes $11\le p\le199$ and at $307$, $353$; above $1$ exactly from $p=17$ on), the bound of Proposition~\ref{prop:wnet}(2) and the explicit pair of (3), which is not a closest pair; open circles: $\mathcal T_p$ (exact, $17\le p\le199$); dotted: the optimum $\mu^*$ of the family (Proposition~\ref{prop:wfamily}), $3.1\cdot10^{-4}$ above $\sqrt{281}/16$.}
